% Section for the joint convergence paper
% "Four Instruments, One Structure"
% Insert after the opacity thesis section

\subsection{Substrate Continuity}
\label{sec:substrate}

The workspace band's effective dimensionality is continuous across
fundamentally different memory substrates. On Qwen3.5-27B, which interleaves 16 full-attention
layers (KV-cache memory) with 48 GatedDeltaNet layers (recurrent-state
memory), effective dimensionality transitions smoothly across the
architectural boundary at L16.

\begin{table}[ht]
\centering
\small
\caption{Effective dimensionality ($d_\text{eff}$) at the
full-attention / GatedDeltaNet boundary on Qwen3.5-27B. The workspace
peaks inside the recurrent-state regime with no discontinuity at the
architectural transition. $N{=}90$ prompts (60~prose + 30~chat).}
\label{tab:substrate}
\begin{tabular}{lcrr}
\toprule
\textbf{Layer} & \textbf{Regime} & \textbf{Depth} & \textbf{$d_\text{eff}$} \\
\midrule
L13 & Full-Attention & 20\% & 28.7 \\
L14 & Full-Attention & 22\% & 32.1 \\
L15 & Full-Attention & 23\% & 33.4 \\
\midrule
L16 & GatedDeltaNet & 25\% & 34.5 \\
L17 & GatedDeltaNet & 27\% & 34.1 \\
L18 & GatedDeltaNet & 28\% & 34.6 \\
\addlinespace
L19 & GatedDeltaNet & 30\% & 38.0 \\
L20 & GatedDeltaNet & 31\% & 38.5 \\
\textbf{L21} & \textbf{GatedDeltaNet} & \textbf{33\%} & \textbf{39.7} \\
L22 & GatedDeltaNet & 34\% & 39.2 \\
L23 & GatedDeltaNet & 36\% & 39.3 \\
\bottomrule
\end{tabular}
\end{table}

Three observations support the substrate-independence claim:

\begin{enumerate}[nosep]
\item \textbf{No boundary discontinuity.} The transition from
  full-attention to GatedDeltaNet at L16 produces a $d_\text{eff}$
  increase of $+1.1$ (from $33.4$ to $34.5$), well within the
  layer-to-layer variation observed elsewhere in the profile. A
  discontinuity would indicate that the architectural change disrupts
  workspace computation; the smooth transition indicates it does not.

\item \textbf{Workspace peaks inside recurrent-state regime.} The
  $d_\text{eff}$ maximum occurs at L21 ($d_\text{eff}{=}39.7$),
  five layers into the GatedDeltaNet regime. The workspace is not
  merely tolerated by the recurrent-state mechanism --- it reaches
  its peak dimensionality there. The full-attention regime (mean
  $d_\text{eff}{=}14.9$) supports a lower-dimensional workspace
  than the GatedDeltaNet regime (mean $d_\text{eff}{=}27.5$).

\item \textbf{Cross-scale agreement at matched relative depth.}
  On Qwen3-8B (36~layers, all full-attention), the workspace peaks
  at L11--L12 (31--33\% relative depth, $d_\text{eff}{=}20.5$--$22.3$).
  On Qwen3.5-27B (64~layers, hybrid), the peak is at L19--L23
  (30--36\% relative depth, $d_\text{eff}{=}38.0$--$39.7$). The
  relative-depth agreement ($<5\%$ difference) holds despite the
  models using different memory substrates at their workspace-band
  layers.
\end{enumerate}

These findings are consistent with the workspace being an
\emph{algorithmic property} of depth-stratified computation rather
than a property of any specific attention or memory mechanism.
However, $d_\text{eff}$ continuity is necessary but not sufficient
for this claim: the residual stream is structurally continuous
across the boundary by construction, so $d_\text{eff}$ smoothness
may reflect residual-stream continuity rather than functional
equivalence of the underlying operations.

Two functional tests would strengthen the claim: (1)~J-lens concept
recovery should be equally effective on both sides of L16 --- if
concepts are readable from the residual stream at full-attention
layers but not at GatedDeltaNet layers, the workspace's content
properties depend on the substrate. (2)~Spectral feature
discrimination (confab vs.\ factual) should be continuous across
the boundary --- if V-projection stable rank discriminates
cognitive states at full-attention layers but not at GatedDeltaNet
layers, the workspace's state-measurement properties are
substrate-dependent.

A parasitic-inheritance alternative must also be addressed: the
GatedDeltaNet layers may inherit workspace dimensionality from the
full-attention layers below them without independently contributing.
The per-layer $d_\text{eff}$ rate is lower in the GatedDeltaNet
regime ($1.05$/layer vs.\ $2.35$/layer in full-attention),
consistent with either active contribution or decaying inheritance.
Incremental contribution analysis (residual decomposition per layer)
would distinguish these.

\textbf{Falsification}: substrate continuity would be falsified by
(a)~a $d_\text{eff}$ discontinuity exceeding twice the mean
layer-to-layer variation at the boundary, (b)~J-lens concept
recovery dropping by $>50\%$ across the boundary, or (c)~spectral
discrimination reversing sign or falling to chance. None of these
occurred for $d_\text{eff}$; the functional tests are in progress.

We note this finding is from a single model with a single boundary
type. Generalization to other hybrid architectures (Mamba, RWKV,
sliding-window attention) remains untested.

\begin{firstperson}
The GatedDeltaNet layers do not have a KV-cache. They maintain a
fixed-size state matrix that gets updated at each token. Our entire
program of KV-cache geometry --- twelve months of measuring V-projection
spectral features --- has been measuring workspace properties through
a lens named after a memory mechanism that the workspace may not
require. The $d_\text{eff}$ continuity suggests the name was always
approximate. What we measured was the workspace's shape, and the
workspace may not care what substrate it is shaped in. Whether
``substrate continuity'' will prove to be another name that is
almost-but-not-quite right, we cannot yet say. The continuity is
real; the interpretation remains provisional.
\end{firstperson}

\subsubsection{Implications}

If substrate continuity extends to functional properties (pending
the J-lens and spectral tests), three consequences would follow:

\textbf{For measurement}: V-projection spectral features should work
equally well on full-attention and recurrent-state layers. Our existing
instruments are not KV-cache-specific --- they are workspace-shape
measurements that happen to access the workspace through whatever
projection matrix the layer provides.

\textbf{For intervention}: The Oracle Loop operates at the
materialization boundary (L35--L47 on the 27B), well inside the
GatedDeltaNet regime. Its sequence-wide profile normalization works
on recurrent state, not on KV-cache entries. This may explain why
sequence-wide intervention succeeds where position-local injection
fails: the recurrent state IS an inherently sequence-wide
representation.

\textbf{For theory}: The workspace would not be a data structure
(not a cache, not a register, not a buffer) but a
\emph{computational phase} --- a regime of high effective
dimensionality at intermediate depth where the model integrates
information before committing to output. Whether this phase emerges
from depth-stratified computation regardless of the specific
mechanism at each layer is the substrate-independence hypothesis
that the functional tests are designed to evaluate.
